\chapter{Convergência TI/TO}
\label{cap:ti-to}

\begin{objetivos}
Ao final deste capítulo você deve ser capaz de:
\begin{itemize}
  \item explicar por que as práticas de TI e de tecnologia de operação (TO) divergem;
  \item dizer o que roda dentro do controlador, o que roda na borda e o que roda fora da fábrica;
  \item avaliar o uso de visão computacional, de modelos de linguagem pequenos, de plugins, de
        Node-RED, de código gerado a partir de modelos e de Linux embarcado;
  \item reconhecer, em cada tecnologia, o ganho, a perda e a exigência que ela traz.
\end{itemize}
\end{objetivos}

\section{Duas culturas, um mesmo sistema}

O chão de fábrica e o data center passaram a compartilhar rede, dados e até hardware. As práticas,
porém, continuam diferentes --- e é essa diferença que explica a maior parte dos atritos de projeto.

\begin{table}[H]
  \centering
  \caption{Tecnologia da informação (TI) e tecnologia de operação (TO), lado a lado.}
  \label{tab:ti-versus-to}
  \begin{tabular}{p{3.0cm} p{5.4cm} p{5.8cm}}
    \toprule
    Aspecto & TI & TO \\
    \midrule
    Prioridade
      & Confidencialidade e integridade do dado
      & Disponibilidade e segurança física do processo \\
    \addlinespace
    Tempo
      & Melhor esforço; a resposta pode variar
      & Prazo garantido, em milissegundos, para as funções de controle \\
    \addlinespace
    Ciclo de vida
      & Atualização frequente, reinício aceitável
      & Equipamento em operação por muitos anos; reinício é parada de produção \\
    \addlinespace
    Mudança
      & Implantação contínua, versão nova a cada semana
      & Mudança planejada, com janela, teste e aprovação \\
    \addlinespace
    Quem opera
      & Equipe de infraestrutura, à distância
      & Equipe de automação e manutenção, no local \\
    \bottomrule
  \end{tabular}
\end{table}

Convergência TI/TO não é escolher uma das colunas da Tabela~\ref{tab:ti-versus-to}. É decidir, item
por item, onde cada prática vale
--- e reconhecer que ``atualizar tudo automaticamente'' e ``prazo garantido de 2\,ms'' são
requisitos incompatíveis no mesmo componente.

\index{TI/TO}\index{tecnologia de operação}

\section{Quem roda onde}

\begin{figure}[H]
  \centering
  \includegraphics[width=0.54\linewidth]{21-ti-to.pdf}
  \caption{Onde cada tecnologia se encaixa, do laço de controle à nuvem.}
  \label{fig:ti-to}
\end{figure}

A Figura~\ref{fig:ti-to} organiza o capítulo inteiro. A regra que ela representa:

\begin{itemize}
  \item no \textbf{controlador} fica o que é determinístico e o que protege a máquina. Só entra ali
        o que tem prazo previsível e comportamento verificável;
  \item na \textbf{borda} --- um computador industrial, um gateway ou um módulo de processamento no
        painel --- ficam as funções que precisam de mais poder de cálculo e admitem latência maior:
        imagem, integração de protocolos, inferência leve;
  \item \textbf{fora da fábrica} ficam o acúmulo de longo prazo e a análise pesada, que não participam
        de nenhum laço de controle.
\end{itemize}

O retorno de informação da borda para o controlador existe, mas passa por um filtro: chega o
\emph{resultado} (uma referência, um valor de correção, um indicador), nunca o modelo nem o
julgamento automático de uma função não determinística.

\section{Visão computacional e processamento de imagem}

Inspeção por imagem deixou de ser exclusividade de grandes linhas. Um conjunto moderno tem quatro
partes: a \textbf{câmera} industrial (padrões de interface como GigE Vision e USB3 Vision existem
justamente para desacoplar câmera de software), a \textbf{iluminação}, a \textbf{óptica} e o
\textbf{processamento}.

Os usos típicos são de verificação, não de medição de precisão:

\begin{itemize}
  \item presença e ausência de componente, tampa, rótulo ou parafuso;
  \item posição e orientação de peça antes do manuseio;
  \item leitura de código de barras e de código bidimensional, com conferência contra a ordem de
        produção;
  \item verificação de cor e de aspecto superficial, com limite de aceitação definido.
\end{itemize}

Para o processamento, a biblioteca aberta OpenCV é a base mais difundida; softwares de inspeção
comerciais montam sobre ela ou sobre bibliotecas próprias.

O ponto de integração com o CLP é mais simples do que parece: o resultado da inspeção é um
\textbf{sinal discreto} --- passa, não passa, peça presente --- entregue por uma saída do sistema de
visão para uma entrada do controlador, ou por variável publicada em OPC UA
(Capítulo~\ref{cap:iiot}). O CLP decide o que fazer com o resultado; não processa a imagem.

\begin{atencao}
Em visão computacional, o resultado depende mais de \emph{iluminação e fixação} do que do algoritmo.
Uma câmera excelente com iluminação instável produz reprovação intermitente, e o operador aprende a
desligar a inspeção. Antes de escolher biblioteca, resolva o gabarito mecânico e a luz --- e defina o
que o sistema deve fazer quando não consegue decidir.
\end{atencao}

\section{Modelos de linguagem pequenos no ambiente de controle}

Modelos de linguagem de pequeno porte (\emph{small language models}, SLM) são modelos compactos, que
podem ser executados localmente, sem depender de serviço externo. É a forma mais realista de levar
capacidade de linguagem para perto do processo, porque mantém o dado dentro da planta.

\paragraph{O que é viável hoje.} Funções de apoio, com o resultado passando por confirmação humana:

\begin{itemize}
  \item \textbf{ajuda ao operador}: responder, em linguagem natural, sobre procedimento, alarme ou
        significado de uma variável, a partir de documentação e do histórico da própria planta;
  \item \textbf{diagnóstico assistido}: receber o conjunto de alarmes e estados no momento da falha e
        sugerir causas prováveis, com base em casos anteriores registrados;
  \item \textbf{geração de relatório}: transformar dados de turno em texto de ocorrência, para revisão
        e assinatura;
  \item \textbf{apoio à engenharia}: explicar um trecho de código, sugerir nome de variável, resumir
        diferenças entre duas revisões de programa.
\end{itemize}

\paragraph{O que é irrealista.} Inferência de modelo dentro do ciclo de varredura. Um modelo, mesmo
pequeno, responde em frações de segundo a alguns segundos e ocupa memória na casa de centenas de
megabytes a alguns gigabytes --- uma a três ordens de grandeza acima do que o ciclo de controle
admite (Capítulo~\ref{cap:funcionamento}). Além do tempo, faltam as propriedades que o controle
exige: comportamento determinístico, resultado reproduzível e verificável, e rastreabilidade da
decisão.

\begin{atencao}
Nenhuma função de segurança, intertravamento ou proteção de máquina pode depender de modelo de
linguagem. Um modelo pode errar com fluência: produz respostas plausíveis e incorretas, sem sinalizar
dúvida. Colocar isso dentro do laço de controle significa aceitar que a máquina tome decisões
imprevisíveis --- e que ninguém consiga explicar depois por que ela tomou. Se a IA vai ajudar, que
ajude \emph{o operador} e \emph{o engenheiro}, com o resultado registrado, conferido e assinado por
gente.
\end{atencao}

Duas questões de governança precisam ser respondidas antes de instalar qualquer assistente:

\begin{enumerate}
  \item \textbf{quem valida a resposta} e como o erro é detectado, antes que vire ação;
  \item \textbf{o que sai da planta}. Enviar dados de processo, receitas e parâmetros para um serviço
        externo é decisão de segurança da informação, não de engenharia de automação
        (Capítulo~\ref{cap:seguranca}).
\end{enumerate}

\section{Plugins e extensões no ambiente de engenharia}

Ambientes de programação aceitam extensões: bibliotecas de blocos, verificadores de padronização,
geradores de documentação e integrações com ferramentas de versionamento. O efeito prático é grande
--- por exemplo, um verificador que aplica automaticamente o checklist do
Capítulo~\ref{cap:boas-praticas} encontra problemas antes da revisão humana.

O que se ganha e o que se perde:

\begin{itemize}
  \item \textbf{ganha}: produtividade, padronização e geração de documentação a partir do próprio
        projeto;
  \item \textbf{perde}: portabilidade. O plugin é do ambiente, não da norma; um projeto que depende
        de extensão proprietária não migra nem para outra versão, em alguns casos;
  \item \textbf{exige}: política de versão. Qual versão do plugin foi usada para gerar aquele
        programa? Sem essa informação, o projeto deixa de ser reconstruível. É o mesmo problema de
        reprodutibilidade que existe em qualquer \emph{build} de software, com a diferença de que
        aqui o artefato controla uma máquina.
\end{itemize}

Há ainda a dimensão de segurança: um plugin executa código dentro do ambiente que fala diretamente
com o controlador. Ele deve ser tratado como software de terceiro em rede industrial --- origem
conhecida, versão controlada e instalado com aprovação.

\index{plugin}\index{ambiente de programação!extensões}

\section{Node-RED}

Node-RED é uma ferramenta de programação visual de fluxos, construída sobre Node.js, muito usada em
borda industrial. Cada nó faz uma coisa --- ler um protocolo, transformar um valor, publicar em um
tópico, enviar um aviso --- e o programador liga os nós.

Os usos em que ela se destaca:

\begin{itemize}
  \item \textbf{prototipagem rápida} de integração, para provar que o dado chega e tem o significado
        esperado;
  \item \textbf{gateway de protocolo}, traduzindo entre Modbus, OPC UA e MQTT
        (Capítulos~\ref{cap:redes-scada} e~\ref{cap:iiot});
  \item \textbf{orquestração de avisos}: consolidar condições e notificar a manutenção;
  \item \textbf{painéis simples} de acompanhamento.
\end{itemize}

E os limites, que precisam estar claros antes de a ferramenta virar infraestrutura crítica:

\begin{enumerate}
  \item \textbf{não é ambiente determinístico}. Roda sobre um sistema de propósito geral, com
        \emph{garbage collector} e rede compartilhada; o tempo de resposta varia;
  \item \textbf{não substitui lógica de segurança}. Nenhum intertravamento pode morar em um fluxo
        Node-RED, ainda que funcione perfeitamente em teste;
  \item \textbf{exige disciplina de manutenção}. Fluxo visual sem versionamento, sem teste e sem
        comentário vira ``caixa preta'' que ninguém se atreve a alterar --- o mesmo problema que a
        documentação do Capítulo~\ref{cap:documentacao} procura evitar;
  \item \textbf{herda os riscos do ambiente}. Portas abertas, credenciais padrão e módulos de
        terceiros desatualizados são o caminho mais comum para um incidente.
\end{enumerate}

A forma segura é tratar Node-RED como parte da camada de \emph{informação}: ele lê, transforma e
publica. Ele não comanda proteção.

\section{Geração de código a partir de modelagem e simulação}

Ambientes de modelagem e simulação --- MATLAB/Simulink é o exemplo mais difundido --- permitem
descrever o comportamento do sistema, simular antes de embarcar e gerar código para o controlador.
O uso típico no curso técnico e na indústria é:

\begin{enumerate}
  \item modelar a planta e o controlador;
  \item testar o algoritmo em simulação, com cenários difíceis, sem risco;
  \item gerar código (texto estruturado ou linguagem de mais baixo nível, conforme o ambiente);
  \item embarcar, testar em bancada e depois em campo.
\end{enumerate}

O ganho é real: algoritmos de controle que seriam penosos de escrever à mão --- filtros, controle
preditivo, escalonamento de ganhos --- passam a ser testáveis antes de tocar a máquina.

Os limites também são reais:

\begin{itemize}
  \item \textbf{código gerado precisa ser revisto}. Ele é correto em relação ao modelo, e o modelo
        pode estar errado em relação à planta;
  \item \textbf{rastreabilidade}: é preciso guardar o modelo, a versão do gerador e o código
        resultante. Alterar o código à mão e perder o vínculo com o modelo é caminho sem volta;
  \item \textbf{nem todo ambiente aceita} o código gerado, ou o aceita com restrições de desempenho
        e de uso de memória;
  \item \textbf{uso em função de segurança} exige processo próprio de verificação, com a norma de
        segurança funcional aplicável (Capítulo~\ref{cap:seguranca}) --- a geração automática não
        substitui a validação.
\end{itemize}

\section{Linux em sistemas embarcados e controladores}
\index{Linux embarcado}

Distribuições Linux viraram base de computadores industriais, gateways de borda, IHMs avançadas e de
controladores que executam \emph{soft-PLC} --- o controlador implementado em software sobre hardware
de propósito geral. As razões são práticas: ecossistema amplo, isolamento por contêineres,
virtualização, atualização gerenciada e integração natural com ferramentas de TI.

O que se encontra no equipamento, porém, raramente é a distribuição de um computador pessoal. O
usual é uma \textbf{imagem sob medida}, gerada pelo fabricante com apenas os componentes
necessários. A ferramenta mais difundida para isso é o \emph{Yocto Project}\index{Yocto}, com a
camada de metadados OpenEmbedded: em vez de instalar um sistema completo e retirar o que sobra,
descrevem-se \emph{receitas} e o resultado é uma imagem mínima, reprodutível e versionada. É o
caminho adotado por vários fabricantes de controladores --- os controladores PLCnext da Phoenix
Contact, por exemplo, executam uma distribuição Linux gerada com o Yocto, conforme a documentação do
fabricante.

A diferença não é só de vocabulário. Uma imagem sob medida chega a campo sem gerenciador de pacotes
aberto ao usuário, sem serviços que ninguém pediu e com uma lista de componentes que se pode
auditar: é menos superfície exposta (Capítulo~\ref{cap:seguranca}) e menos coisa para quebrar numa
atualização. Em troca, acrescentar um \emph{driver} ou um programa novo deixa de ser um comando e
passa a depender da cadeia de compilação cruzada do fabricante --- negocia-se escopo, versão e
prazo, e não se resolve na madrugada anterior à partida. Vale saber sobre que base roda o seu
controlador antes de prometer uma integração que exija pacote de terceiro.

O ponto que interessa à automação é o \textbf{determinismo}. Linux comum prioriza aproveitamento e
não prazo: a resposta pode variar conforme a carga. Núcleos com preempção de tempo real reduzem
bastante essa variabilidade e são adequados a muitas tarefas de controle. Ainda assim, o pior caso
de um sistema de propósito geral com \emph{soft-PLC} não é o mesmo de um controlador dedicado com
sistema operacional de tempo real, e a diferença está justamente no que a norma de segurança
funcional exige demonstrar.

\begin{table}[H]
  \centering
  \caption{Usos de Linux na automação e o cuidado correspondente.}
  \label{tab:linux-uso}
  \begin{tabular}{p{4.6cm} p{4.2cm} p{5.5cm}}
    \toprule
    Uso & É adequado? & Cuidado \\
    \midrule
    Gateway de borda, integração e histórico local
      & Sim, é o uso mais comum
      & Manter atualizado e sem serviços desnecessários expostos \\
    \addlinespace
    IHM e supervisão em computador industrial
      & Sim
      & Isolar de outras funções; não acumular servidor de arquivos e supervisão na mesma máquina \\
    \addlinespace
    \emph{Soft-PLC} em controle de máquina
      & Depende do requisito de prazo
      & Medir o pior caso e verificar contra o requisito do processo, não contra a média \\
    \addlinespace
    Função de segurança
      & Só com hardware e sistema dedicados e certificados
      & Não usar computador de propósito geral como elemento de segurança \\
    \bottomrule
  \end{tabular}
\end{table}

\section{O balanço: o que ganha, o que perde, o que exige}

\begin{table}[H]
  \centering
  \caption{Balanço de cada tecnologia de TI trazida para a operação.}
  \label{tab:balanco-ti-to}
  \begin{tabular}{p{2.9cm} p{3.9cm} p{3.4cm} p{4.1cm}}
    \toprule
    Tecnologia & O que ganha & O que perde & O que exige \\
    \midrule
    Visão computacional
      & Inspeção automática de características difíceis de medir por sensor
      & Determinismo do resultado, se a iluminação ou a peça variar
      & Gabarito mecânico, iluminação estável e critério de decisão escrito \\
    \addlinespace
    SLM na instalação
      & Apoio ao operador e ao engenheiro; dado não sai da planta
      & Confiabilidade da resposta; nenhuma garantia de prazo
      & Ficar fora do laço de controle, com resultado conferido e registrado \\
    \addlinespace
    Plugin no IDE
      & Produtividade, padronização e documentação automática
      & Portabilidade do projeto
      & Política de versão e aprovação de software de terceiro \\
    \addlinespace
    Node-RED
      & Integração rápida entre protocolos e prototipagem
      & Determinismo e rastreabilidade
      & Versionar os fluxos e nunca usá-lo para proteção \\
    \addlinespace
    Código gerado por modelo
      & Algoritmo complexo testado em simulação antes de embarcar
      & Vínculo com o modelo, se o código for editado à mão
      & Rastreabilidade entre modelo, gerador e código embarcado \\
    \addlinespace
    Linux embarcado
      & Flexibilidade, atualização e integração com TI
      & Pior caso de tempo, frente a um controlador dedicado
      & Medir o pior caso e mantê-lo atualizado \\
    \bottomrule
  \end{tabular}
\end{table}

Nenhuma linha da Tabela~\ref{tab:balanco-ti-to} é gratuita: cada tecnologia se paga em um contexto e
cobra em outro, e a decisão depende do que a máquina precisa garantir --- prazo, segurança ou
simplesmente sair do papel.

\begin{pratica}
Escolha uma das tecnologias da Tabela~\ref{tab:balanco-ti-to} e prepare uma proposta de uso em uma
máquina real, respondendo:
(a) qual problema concreto ela resolve, medido em indicador?
(b) em qual camada da Figura~\ref{fig:ti-to} ela vai rodar, e por quê?
(c) qual é o pior caso de tempo aceitável, e como você o mediria?
(d) o que acontece com a máquina se essa tecnologia falhar ou parar de responder?
(e) quem mantém isso depois de você sair?
\end{pratica}

\begin{exercicios}
\begin{enumerate}[label=\alph*)]
  \item Explique, em três linhas, por que atualização automática e prazo garantido são requisitos
        incompatíveis no mesmo componente.
  \item Localize na Figura~\ref{fig:ti-to} onde roda um verificador de padrões de código e onde roda
        um modelo de linguagem que resume alarmes.
  \item Cite quatro usos típicos de visão computacional em linha de montagem e diga qual deles é o
        mais fácil de implantar.
  \item Por que, em visão computacional, iluminação e fixação importam mais que o algoritmo?
  \item Como o resultado de uma inspeção por imagem chega ao CLP? Descreva dois caminhos.
  \item Liste três usos viáveis de SLM em uma planta e explique, para cada um, onde entra a
        confirmação humana.
  \item Explique por que a inferência de modelo não cabe no ciclo de varredura, citando tempo e
        memória.
  \item Um integrador propõe usar um assistente de IA para decidir o bloqueio de uma prensa.
        Escreva a recusa técnica, com dois argumentos.
  \item Cite dois ganhos e duas perdas de adotar plugin de terceiro no ambiente de engenharia.
  \item Por que Node-RED não pode executar intertravamento, mesmo que o teste tenha funcionado?
  \item Descreva o problema de rastreabilidade criado quando alguém edita à mão um código gerado
        automaticamente.
  \item Aplique a Tabela~\ref{tab:linux-uso} a um computador industrial que já executa supervisão e
        passa a executar também \emph{soft-PLC} de uma máquina. Justifique sua decisão.
  \item Escolha uma linha da Tabela~\ref{tab:balanco-ti-to} e escreva a exigência técnica que o
        ``o que exige'' representa, em forma de cláusula de especificação.
\end{enumerate}
\end{exercicios}
